\chapter{کدنویسی و تحقیق با منبع}

ابزار فایل متن 
\lr{UTF-8}
 و کد کوچک تا ۶۴ کیلوبایت را می‌پذیرد. مسیر مطلق، رفتن به پوشهٔ والد، خروج از پوشهٔ کاری، اجزای مخفی، فایل‌های حافظه، پوشه‌های وابستگی و نام‌های رایج اطلاعات محرمانه مسدودند. فهرست پروژه حداکثر ۲۰۰ فایل مجاز نشان می‌دهد. متن فایل دادهٔ ورودی است؛ توضیح داخل کد به دستور سیستمی تبدیل نمی‌شود.

ابتدا ابزار مستقیم را امتحان کنید تا بدانید مدل چه چیزی می‌بیند:

\begin{LTR}
\begin{Verbatim}[breaklines,breakanywhere,fontsize=\small]
/files
/read README.md
\end{Verbatim}
\end{LTR}

برای تمرین کدنویسی، در پوشهٔ نمونهٔ جداگانه فایل 
\lr{demo.py}
 را بسازید:

\begin{LTR}
\begin{Verbatim}[breaklines,breakanywhere,fontsize=\small]
def add(a, b):
    return a - b
\end{Verbatim}
\end{LTR}

سپس بنویسید: «فایل 
\lr{demo.py}
 را بخوان. توضیح بده چرا 
\lr{add}
 اشتباه است. کوچک‌ترین اصلاح را پیشنهاد کن و فقط بعد از نشان دادن تفاوت ذخیره کن». دستیار باید 
\lr{read\_file}
 را اجرا کند، تفریق را تشخیص دهد و با 
\lr{write\_file}
 متن کامل اصلاح‌شده را پیشنهاد دهد. مسیر و تفاوت را بررسی کنید. 
\lr{YES}
 تغییر را ذخیره و 
\lr{Enter}
 لغو می‌کند. ابزار بررسی می‌کند فایل در زمان تأیید شما توسط فرد دیگری تغییر نکرده باشد. ذخیرهٔ موفق یعنی فایل نوشته شده؛ درستی برنامه را ثابت نمی‌کند.

اصلاح کوچک مورد انتظار:

\begin{LTR}
\begin{Verbatim}[breaklines,breakanywhere,fontsize=\small]
def add(a, b):
    return a + b
\end{Verbatim}
\end{LTR}

در همان پوشه‌ای که 
\lr{demo.py}
 قرار دارد، خودتان بررسی واقعی را اجرا کنید. ابزار اجرای کد در ایجنت نیست و نباید ادعا کند این دستور را قبلاً اجرا کرده است:

\begin{LTR}
\begin{Verbatim}[breaklines,breakanywhere,fontsize=\small]
python -c "from demo import add; assert add(2, 3) == 5; print('PASS')"
git diff
\end{Verbatim}
\end{LTR}

جایگزینی کامل فایل برای یادگیری ساده است، اما برای فایل کوچک طراحی شده است. تبدیل 
\lr{AST}
، پچ معنایی، جستجوی مخزن بزرگ، 
\lr{commit}
 خودکار و سامانهٔ 
\lr{rollback}
 وجود ندارد. قبل از 
\lr{commit}
 تفاوت 
\lr{Git}
 را ببینید. این نسخه مخزن پروژهٔ دیگر را از طرف شما ایجاد یا 
\lr{commit}
 نمی‌کند.

تحقیق دو مرحله دارد: سرچ متن کوتاه و 
\lr{URL}
 می‌دهد؛ خواندن صفحه متن 
\lr{HTML}
 را دریافت می‌کند. درخواست خوب هر دو را می‌خواهد و مشخص می‌کند کدام ادعا در متن منبع پشتیبانی می‌شود.

\begin{LTR}
\begin{Verbatim}[breaklines,breakanywhere,fontsize=\small]
/search Python official tutorial list comprehensions
/open https://docs.python.org/3/tutorial/datastructures.html
\end{Verbatim}
\end{LTR}

نمونه: «مستندات رسمی 
\lr{Python}
 دربارهٔ 
\lr{list comprehension}
 را پیدا کن، صفحهٔ مرتبط را بخوان و سه نکته را با 
\lr{URL}
 منبع توضیح بده. اگر جزئیاتی در صفحه نیست، صریح بگو». عنوان و نشانی نتیجهٔ سرچ فقط سرنخ است. متن صفحه ممکن است کوتاه‌شده، قدیمی، ناقص یا اشتباه باشد. نقل‌قول باید عیناً در متن دریافت‌شده وجود داشته باشد؛ در غیر این صورت باید بازگویی باشد.

خوانندهٔ وب فقط 
\lr{HTTPS}
 عمومی روی پورت ۴۴۳ می‌پذیرد، مقصد هر تغییر مسیر را دوباره بررسی می‌کند، نشانی خصوصی یا رزروشده را رد می‌کند و به همان 
\lr{IP}
 بررسی‌شده با کنترل نام میزبان در 
\lr{TLS}
 وصل می‌شود. محتوا بدون فشرده‌سازی درخواست می‌شود، سقف دریافت ۲ مگابایت و سقف متن خروجی ۱۲ هزار نویسه است. ورود به سایت، اجرای 
\lr{JavaScript}
 و خواندن 
\lr{PDF}
 ندارد. صفحهٔ ورود یا محافظت ربات ممکن است شواهد ناقص یا غیرقابل‌دریافت بدهد. شکست سرچ به‌صورت خطا گزارش می‌شود، نه نتیجهٔ ساختگی.

مقابله با دستورهای مخرب دو لایه دارد: پرامپت محتوای وب و فایل و حافظه را نامطمئن معرفی می‌کند؛ کد ابزار نیز مسیر، آرگومان، نشانی سرویس و تأیید تغییر را محدود می‌کند. پرامپت به‌تنهایی تضمین نمی‌کند مدل هیچ‌گاه متن مخرب را دنبال نکند. تغییر و شواهد را همچنان بررسی کنید. مدت 
\lr{DNS}
 به 
\lr{resolver}
 سیستم‌عامل وابسته است؛ بودجهٔ زمان خوانندهٔ وب محدودیت عملی است و تضمین مطلق برای همهٔ عملیات 
\lr{resolver}
 نیست.

\cite{ddgs}
